\subsubsection{Trao đổi khóa AES qua kênh mã hóa RSA}
\indent Sau khi hoàn tất quá trình mã hóa nội dung tin nhắn bằng \textbf{AES}, ứng dụng cần gửi khóa AES này cho những người nhận trong cuộc trò chuyện để họ có thể tiến hành giải mã.  
Tuy nhiên, nếu gửi trực tiếp khóa AES thì sẽ tiềm ẩn rủi ro lộ thông tin. Vì vậy, hệ thống áp dụng cơ chế mã hóa bất đối xứng \textbf{RSA} để bảo vệ khóa phiên trước khi truyền đi.

\indent Trước khi gửi tin nhắn, ứng dụng gửi yêu cầu đến máy chủ để lấy danh sách thành viên của cuộc trò chuyện, bao gồm khóa công khai \textbf{RSA} của từng nhận.  
Với khóa công khai này, ứng dụng tiến hành mã hóa khóa AES phiên riêng biệt cho từng thành viên.  
Cụ thể, cùng một khóa AES sẽ được mã hóa nhiều lần, mỗi lần bằng khóa công khai của một người nhận khác nhau.

\indent Kết quả là mỗi người dùng sẽ nhận được một phiên bản khóa AES được mã hóa dành riêng cho họ, đảm bảo rằng chỉ thiết bị của họ mới có thể giải mã bằng khóa bí mật \textbf{RSA} lưu trữ cục bộ.  
Danh sách các khóa AES đã mã hóa được gửi kèm trong gói tin nhắn lên máy chủ, dưới dạng cặp \texttt{\{userId, encryptedAesKey\}}.

\indent Trong toàn bộ quy trình này, máy chủ chỉ đóng vai trò trung gian truyền dữ liệu, hoàn toàn không có khả năng giải mã bất kỳ khóa AES nào vì không sở hữu khóa bí mật tương ứng.  
Cách triển khai này giúp đảm bảo khóa phiên chỉ có thể được giải mã tại phía người nhận, duy trì tính bảo mật tuyệt đối cho mô hình mã hóa đầu cuối (\textbf{E2EE}).
